\documentclass[../main.tex]{subfiles}

\begin{document}

\chapter{Elements of the document}
\label{chap:tutorial3}

\section{Images}

\begin{figure}[H]
\centering
\includegraphics[width=0.5\textwidth]{u14.jpg}
\caption{Example image}
\label{img:ref_image}
\end{figure}

Reference to the image \ref{img:ref_image}.

\section{Tables}
\begin{table}
\centering
\begin{tabular}{|c|c|c|} 
 \hline
 cell1 & cell2 & cell3 \\ \hline
 cell4 & cell5 & cell6 \\ \hline
 cell7 & cell8 & cell9 \\ \hline
\end{tabular}
\caption{Example table}
\label{tab:ref_table}
\end{table}

Reference to the table \ref{tab:ref_table}.

\section{List}
Example of an unordered list:
\begin{itemize}
  \item List entries start with the \verb|\item| command.
  \item Individual entries are indicated with a black dot, a so-called bullet.
  \item The text in the entries may be of any length.
\end{itemize}

Example of an ordered list:
\begin{enumerate}
  \item This is the first entry in our list.
  \item The list numbers increase with each entry we add.
\end{enumerate}

Example of nested list:
\begin{enumerate}
   \item First level item
   \item First level item
   \begin{enumerate}
     \item Second level item
     \item Second level item
     \begin{enumerate}
       \item Third level item
       \item Third level item
       \begin{enumerate}
         \item Fourth level item
         \item Fourth level item
       \end{enumerate}
     \end{enumerate}
   \end{enumerate}
 \end{enumerate}


\section{Equations}

This is an example of an in-line equation: $E = MC^2$\\
This is an example of an display mode equation:
\begin{equation}
    E = MC^2
\label{eq}
\end{equation}

This is a reference to the Equation \ref{eq}.

\section{Code Snippet}

\subsection{verbatim}
\begin{verbatim}
Text enclosed inside \texttt{verbatim} environment 
is printed directly 
and all \LaTeX{} commands are ignored.
\end{verbatim}

\subsection{listings}
\begin{lstlisting}
import numpy as np
    
def incmatrix(genl1,genl2):
    m = len(genl1)
    n = len(genl2)
    M = None #to become the incidence matrix
    VT = np.zeros((n*m,1), int)  #dummy variable
    
    #compute the bitwise xor matrix
    M1 = bitxormatrix(genl1)
    M2 = np.triu(bitxormatrix(genl2),1) 

    for i in range(m-1):
        for j in range(i+1, m):
            [r,c] = np.where(M2 == M1[i,j])
            for k in range(len(r)):
                VT[(i)*n + r[k]] = 1;
                VT[(i)*n + c[k]] = 1;
                VT[(j)*n + r[k]] = 1;
                VT[(j)*n + c[k]] = 1;
                
                if M is None:
                    M = np.copy(VT)
                else:
                    M = np.concatenate((M, VT), 1)
                
                VT = np.zeros((n*m,1), int)
    
    return M
\end{lstlisting}

\subsection{algpseudocode}
\begin{algorithmic}
\State $i \gets 10$
\If{$i\geq 5$} 
    \State $i \gets i-1$
\Else
    \If{$i\leq 3$}
        \State $i \gets i+2$
    \EndIf
\EndIf 
\end{algorithmic}

\subsection{algorithm}

reference to algorithm \ref{alg:cap}
\begin{algorithm}[!h]
\caption{An algorithm with caption}\label{alg:cap}
\begin{algorithmic}
\Require $n \geq 0$
\Ensure $y = x^n$
\State $y \gets 1$
\State $X \gets x$
\State $N \gets n$
\While{$N \neq 0$}
\If{$N$ is even}
    \State $X \gets X \times X$
    \State $N \gets \frac{N}{2}$  \Comment{This is a comment}
\ElsIf{$N$ is odd}
    \State $y \gets y \times X$
    \State $N \gets N - 1$
\EndIf
\EndWhile
\end{algorithmic}
\end{algorithm}

\section{href}
For further references see \href{http://www.overleaf.com}{Something Linky} 
or go to the next url: \url{http://www.overleaf.com}

\end{document}